% Generated draft: requires mathematical and source review.
\documentclass[11pt]{article}
\usepackage[margin=25mm]{geometry}
\usepackage{fontspec}
\setmainfont{Latin Modern Roman}
\usepackage{amsmath,amssymb,mathtools,unicode-math}
\setmathfont{Latin Modern Math}
\usepackage{enumitem,needspace,xurl,hyperref}
\setlength{\parindent}{0pt}
\setlength{\parskip}{5pt}
\setlength{\emergencystretch}{3em}
\newcommand{\N}{\mathbb{N}}
\newcommand{\Z}{\mathbb{Z}}
\newcommand{\Q}{\mathbb{Q}}
\newcommand{\R}{\mathbb{R}}
\newcommand{\C}{\mathbb{C}}
\newcommand{\F}{\mathbb{F}}
\newcommand{\A}{\mathbb{A}}
\newcommand{\PP}{\mathbb{P}}
\newcommand{\E}{\mathbb{E}}
\newcommand{\eps}{\varepsilon}
\newcommand{\dd}{\,\mathrm{d}}
\begin{document}
\section*{UnsolvedMath ID: 30006100}
\textbf{Source problem code:} catalogue entry
\subsection{Definitions and mathematical statement}
The supplied record asks for quantitative comparison of an exact auxiliary-function optimization problem with its extended dynamic mode decomposition (EDMD) approximation. Its transcription omits problem (1) and does not specify the state space, admissible auxiliary functions, or sampling assumptions. The following notation therefore gives a conditional formulation of the requested comparison, with these omissions left explicit.

Let \(X\subseteq\R^q\), where the dimension \(q\) and the hypotheses on \(X\) remain unspecified. Let \(f\) be a vector field and \(g:X\to\R\) an observable. For an admissible differentiable auxiliary function \(v\), write
\[
\mathcal L v(x)=f(x)\cdot\nabla v(x).
\]
If \(\mathcal V\) denotes the admissible class intended in the missing problem (1), its exact auxiliary-function optimum has the formal expression
\[
U_{\mathrm{ex}}(\mathcal V)
:=\inf\bigl\{U\in\R:\ \exists v\in\mathcal V,\quad
U-g(x)-\mathcal L v(x)\geq0\quad\forall x\in X\bigr\}.
\]
The supplied excerpt does not determine \(\mathcal V\), the hypotheses ensuring that this value is finite, or the hypotheses relating it to a maximal time average.

Let \(\mathcal P_k\) denote the polynomials in \(q\) variables of total degree at most \(k\). Interpreting the excerpt's polynomial degree convention this way, let \(b(x)\in\R^{r_d}\) and \(\Psi(x)\in\R^{r_m}\) be basis vectors for \(\mathcal P_d\) and \(\mathcal P_m\), respectively, where \(m\geq d\geq0\) are integers and \(r_k=\dim\mathcal P_k\). Let \(N\geq1\) be an integer, \(\tau>0\) a sampling step, and \((x_n,y_n)_{n=1}^N\) the supplied type of paired observations. For a flow \(\Phi^t\) generated by \(f\), the intended noise-free interpretation is \(y_n=\Phi^\tau(x_n)\), and the Koopman operator is \((\mathcal K^\tau h)(x)=h(\Phi^\tau(x))\), wherever these expressions exist. Whether the actual observations satisfy this relation, and whether noise is allowed, are unspecified.

Choose a least-squares minimizer
\[
K_{m,N}^{\tau}\in\operatorname*{arg\,min}_{K\in\R^{r_d\times r_m}}
\frac1N\sum_{n=1}^N\left\|b(y_n)-K\Psi(x_n)\right\|_2^2.
\]
This reconstructs the transcribed approximation \(\mathcal K^\tau b\approx K_{m,N}^{\tau}\Psi\). For \(v(x)=c^{\mathsf T}b(x)\in\mathcal P_d\), define
\[
\mathcal L_{m,N}^{\tau}v(x)
:=\frac{c^{\mathsf T}\bigl(K_{m,N}^{\tau}\Psi(x)-b(x)\bigr)}{\tau}.
\]
The corresponding approximate optimum in the supplied problem (2) is
\[
\widehat U_{d,m,N}^{\tau}
:=\inf\bigl\{U\in\R:\ \exists c\in\R^{r_d},\quad
U-g(x)-\mathcal L_{m,N}^{\tau}(c^{\mathsf T}b)(x)\geq0
\quad\forall x\in X\bigr\}.
\]
No uniqueness or selection rule for \(K_{m,N}^{\tau}\) is supplied; the approximate optimum can depend on this choice.

\textbf{Requested statement.} Derive explicit estimates for the difference between the optimal values of the source's problems (2) and (1), displaying their dependence on \(N\), \(m\), and \(\tau\). For finite optimal values, an absolute-gap version of this request would have the form
\[
\left|\widehat U_{d,m,N}^{\tau}-U_{\mathrm{ex}}(\mathcal V)\right|
\leq E(N,m,\tau;\text{specified system, dictionary, and sampling parameters}).
\]
The source does not specify an absolute versus one-sided estimate, a deterministic versus probabilistic guarantee, or a convergence regime. If problem (1) restricts its auxiliary functions to \(\mathcal P_d\), then \(\mathcal V=\mathcal P_d\) gives the corresponding comparison; if it allows a larger class, the comparison also involves the restriction of auxiliary functions to \(\mathcal P_d\). The omitted problem (1) prevents choosing between these targets. Any precise bound must state the additional hypotheses and parameter dependence it requires. The source describes the approximate values as approximate upper bounds for a maximal time average; the supplied material does not establish that they are rigorous upper bounds for the original dynamics.
\subsection{Short English statement}
Estimate how the optimal value of an auxiliary-function problem changes when its exact Lie derivative is replaced by an EDMD approximation, with explicit dependence on the data count \(N\), dictionary degree \(m\), and sampling step \(\tau\). The exact admissible class, dynamical assumptions, sampling model, and type of error guarantee are missing, so this is a conditional statement-only draft requiring human mathematical review.
\subsection{Sources}
\begin{itemize}
\item Local UnsolvedMath registry, supplied complete record for permanent ID \(30006100\), title: \emph{Finite-Sample Error Bounds for EDMD Time-Average Optimization}; exact source problem code: \texttt{catalogue entry}. The record supplies the status \texttt{open} and statement status \texttt{corrected\_verified}; these are registry labels, not independent mathematical verification.
\item Supplied source citation: \emph{Mini-Workshop: Data-driven Modeling, Analysis, and Control of Dynamical Systems} (2025), Oberwolfach Reports, pp.\ 3255--3300, DOI \texttt{10.4171/owr/2024/57}. Supplied source URL: \url{https://doi.org/10.4171/owr/2024/57}. The background describes Oberwolfach Reports as documenting workshops and their open-problem discussions. The record gives proposed year 2024, while its citation gives 2025; the date discrepancy is unresolved.
\item Supplied literature source and attribution: Fantuzzi and Bramburger in Oberwolfach Report 57, \url{https://doi.org/10.4171/OWR/2024/57}. No individual proposer is specified in the registry record.
\item Supplied literature source: Llamazares-Elias et al., EDMD convergence bounds, \url{https://arxiv.org/abs/2405.00539}. The excerpt asks whether recent quantitative EDMD error analysis can help obtain the requested optimal-value estimates.
\item Literature assessment dated exactly \texttt{2026-08-22}: the supplied curation labels the problem open and states that existing EDMD bounds do not provide the complete requested dependence on data count, dictionary degree, and sampling step. This is supplied dated curation, not an independently verified literature assessment. The background explicitly requests independent verification of time-sensitive status and cited claims. No source was visited for this draft.
\end{itemize}
\end{document}
